\paragraph{Larger-bound development screen.}\label{app:higherbound}
The completed screen uses seeds 21, 22 and 23 and the same 400-image
development set. Each trial raises both upper bounds, retains the original
lower bounds, and reuses its saved four-direction mask warm-up and fixed
training readouts. All trials perform the same 2,500 plus 500 guided
updates with paired image sampling and noise. Saved controls are reused.
Table~\ref{tab:higherbound} reports means across the training starts. These
scores measure development performance, not held-out confirmation.
Group membership remains fixed.

\begin{table*}[t]\centering\small
\caption{Values with $\pm$ show mean and sample SD. Seed development means (\%). Each cap passes the original
native-reference MSE limit in all runs.}\label{tab:higherbound}
\begin{tabular}{@{}lrrr@{}}\hline
Upper bound & Capture & Patch AP & VE\\\hline
0.20 (saved) & $37.52\pm0.14$ & $80.10\pm0.39$ & $65.35\pm0.06$\\
0.40 & $38.55\pm0.13$ & $80.68\pm0.35$ & $64.93\pm0.07$\\
0.60 & $38.56\pm0.12$ & $80.68\pm0.35$ & $64.92\pm0.08$\\\hline
\end{tabular}\end{table*}

Raising the cap to 0.40 improves Capture and AP but reduces reconstruction
VE. Increasing it to 0.60 adds only 0.014 Capture percentage points, with
essentially unchanged AP and a 0.005-point VE decrease. These development
results motivated cap 0.40 for the subsequent test audit.

Larger caps do not force larger weights. At the 0.60 endpoints, contrast
weights range from 0.340 to 0.381 and detection weights from 0.273 to
0.327, below the cap. In sampled gradient diagnostics at updates 250 and
500 of the final phase, the cosine between reconstruction and contrast
encoder gradients ranges from $-0.147$ to $-0.088$ across seeds. Detection
gradients have positive cosine with both other losses at these checks.
Meanwhile, the last 500 updates still add 0.50--0.80 Capture points and
0.28--0.46 AP points, while losing 0.11--0.14 VE points. Thus, the cap
plateau does not establish that training has plateaued. These diagnostics
motivate testing gradient conflict as a source of the reconstruction
tradeoff. They do not establish its cause or prove an intervention works.

Complete per-seed development scores, endpoint identities and final-phase
histories are retained in \path{evidence/cap_bounds/}. The original
seed-21-only screen remains in \path{evidence/weight_bound/}. The 0.40 cap
was selected before its test audit. 0.60 was not evaluated on test.

